\documentclass[10pt,a4paper]{amsart}
\usepackage{amsmath,amssymb,amsthm,graphicx}
\usepackage{verbatim}
\newenvironment{reference}{\comment}{\endcomment}
%\newenvironment{reference}{}{}
\newenvironment{slogan}{\comment}{\endcomment}
\newenvironment{history}{\comment}{\endcomment}
% For commutative diagrams we use Xy-pic
\usepackage[all]{xy}

% We use 2cell for 2-commutative diagrams.
\xyoption{2cell}
% Theorem environments.
%
\theoremstyle{plain}
\newtheorem{theorem}[subsection]{Theorem}
\newtheorem{proposition}[subsection]{Proposition}
\newtheorem{lemma}[subsection]{Lemma}

\theoremstyle{definition}
\newtheorem{definition}[subsection]{Definition}
\newtheorem{example}[subsection]{Example}
\newtheorem{exercise}[subsection]{Exercise}
\newtheorem{situation}[subsection]{Situation}

\theoremstyle{remark}
\newtheorem{remark}[subsection]{Remark}
\newtheorem{remarks}[subsection]{Remarks}

\numberwithin{equation}{subsection}

% Macros
%
\def\lim{\mathop{\mathrm{lim}}\nolimits}
\def\colim{\mathop{\mathrm{colim}}\nolimits}
\def\Spec{\mathop{\mathrm{Spec}}}
\def\Hom{\mathop{\mathrm{Hom}}\nolimits}
\def\Ext{\mathop{\mathrm{Ext}}\nolimits}
\def\SheafHom{\mathop{\mathcal{H}\!\mathit{om}}\nolimits}
\def\SheafExt{\mathop{\mathcal{E}\!\mathit{xt}}\nolimits}
\def\Sch{\mathit{Sch}}
\def\Mor{\mathop{\mathrm{Mor}}\nolimits}
\def\Ob{\mathop{\mathrm{Ob}}\nolimits}
\def\Sh{\mathop{\mathit{Sh}}\nolimits}
\def\NL{\mathop{N\!L}\nolimits}
\def\CH{\mathop{\mathrm{CH}}\nolimits}
\def\proetale{{pro\text{-}\acute{e}tale}}
\def\etale{{\acute{e}tale}}
\def\QCoh{\mathit{QCoh}}
\def\Ker{\mathop{\mathrm{Ker}}}
\def\Im{\mathop{\mathrm{Im}}}
\def\Coker{\mathop{\mathrm{Coker}}}
\def\Coim{\mathop{\mathrm{Coim}}}

% Boxtimes
%
\DeclareMathSymbol{\boxtimes}{\mathbin}{AMSa}{"02}

%
% Macros for moduli stacks/spaces
%
\def\QCohstack{\mathcal{QC}\!\mathit{oh}}
\def\Cohstack{\mathcal{C}\!\mathit{oh}}
\def\Spacesstack{\mathcal{S}\!\mathit{paces}}
\def\Quotfunctor{\mathrm{Quot}}
\def\Hilbfunctor{\mathrm{Hilb}}
\def\Curvesstack{\mathcal{C}\!\mathit{urves}}
\def\Polarizedstack{\mathcal{P}\!\mathit{olarized}}
\def\Complexesstack{\mathcal{C}\!\mathit{omplexes}}
% \Pic is the operator that assigns to X its picard group, usage \Pic(X)
% \Picardstack_{X/B} denotes the Picard stack of X over B
% \Picardfunctor_{X/B} denotes the Picard functor of X over B
\def\Pic{\mathop{\mathrm{Pic}}\nolimits}
\def\Picardstack{\mathcal{P}\!\mathit{ic}}
\def\Picardfunctor{\mathrm{Pic}}
\def\Deformationcategory{\mathcal{D}\!\mathit{ef}}
\usepackage{cleveref}
\crefname{originallemma}{original lemma}{original lemmas}
\Crefname{originallemma}{Original lemma}{Original lemmas}
\makeatletter
\newlabel{lemma-colimits}{{073E}{source}{}{sourcepointer-1}{}}
\newlabel{lemma-colimits@cref}{{[originallemma][0][]073E}{[1][0][]source}{}{}{}}
\makeatother
\crefname{originalreference}{original reference}{original references}
\Crefname{originalreference}{Original reference}{Original references}
\makeatletter
\newlabel{situation-vanishing}{{07UY}{source}{}{sourcepointer-2}{}}
\newlabel{situation-vanishing@cref}{{[originalreference][0][]07UY}{[1][0][]source}{}{}{}}
\makeatother
\crefname{originallemma}{original lemma}{original lemmas}
\Crefname{originallemma}{Original lemma}{Original lemmas}
\makeatletter
\newlabel{spaces-morphisms-lemma-pushforward}{{03M9}{source}{}{sourcepointer-3}{}}
\newlabel{spaces-morphisms-lemma-pushforward@cref}{{[originallemma][0][]03M9}{[1][0][]source}{}{}{}}
\makeatother
\title{The pushforward and the global sections of the pullback of a module under the structure morphism in the vanishing situation}
\author{The Stacks Project Authors}
\date{stacks-project:spaces-cohomology.tex:local-snapshot-20261012; GNU FDL 1.2 or later}
\begin{document}
\maketitle
\noindent\textit{Editorial note.} All delivered mathematical text is the author's own, copied from the pinned Stacks Project chapter file spaces-cohomology.tex; no statement or proof step was rewritten, reordered or removed.\par\medskip
\noindent\textit{Editorial note.} Preamble provenance: the pinned chapter file opens with \textbackslash{}input\{preamble\} and carries no preamble text of its own, so the project's single Stacks-wide preamble staged at sources/stacks-spaces-cohomology/preamble.tex is excerpted verbatim as three spans, lines 8-12 (verbatim package and the project's own reference/slogan/history display environments), lines 14-18 (Xy-pic setup) and lines 154-217 (theorem declarations and mathematical macros). The wrapper class and the amsmath/amssymb/amsthm/graphicx base packages are supplied by the assembler. No macro was invented or redefined.\par\medskip
\noindent\textit{Editorial note.} Same-chapter pointers are located at their \textbackslash{}label lines inside the pinned chapter file; cross-chapter pointers are located in the project's own tag list staged in the snapshot (sources/stacks-chow/tags.txt), at the verbatim lines carrying their canonical tags.\par\medskip
\noindent\textit{Allowed standard prerequisites (editorial).}\begin{itemize}
\item Quasi-coherence of the pushforward of a quasi-coherent module. Morphisms of Spaces, label spaces-morphisms-lemma-pushforward (invoked at proof line 3019). It shows that \$p\_*(M \textbackslash{}otimes\_A \textbackslash{}mathcal\{O\}\_X)\$ is quasi-coherent on \$\textbackslash{}Spec(A)\$, so that the equality of global sections determines it as \$\textbackslash{}widetilde\{M\}\$.
\item Global sections of direct sums of the structure sheaf. Same chapter spaces-cohomology.tex, label lemma-colimits (invoked at proof line 3022). It gives \$\textbackslash{}Gamma(X, \textbackslash{}bigoplus\_\{i \textbackslash{}in I\}\textbackslash{}mathcal\{O\}\_X) = \textbackslash{}bigoplus\_\{i \textbackslash{}in I\}A\$, which is the free case of the claim.
\item Vanishing of \$H\textasciicircum{}1\$ in the situation, which makes global sections right exact. Same chapter spaces-cohomology.tex, Situation \textbackslash{}ref\{situation-vanishing\} and the vanishing statement of the surrounding section (used at proof line 3025). It lets the global sections of a presentation of \$M\$ compute the global sections of \$M \textbackslash{}otimes\_A \textbackslash{}mathcal\{O\}\_X\$.
\end{itemize}\medskip
\ifdefined\hypertarget\hypertarget{sourcepointer-1}{}\fi
\noindent\textit{Source reference.} Original lemma 073E: sources/stacks-spaces-cohomology/pinned-original-source.tex, source lines 667--667. Navigation pointer for \textbackslash{}ref\{lemma-colimits\} used inside the counted span (line 3022); the argument that consumes it is retained in the span.\par\medskip
\ifdefined\hypertarget\hypertarget{sourcepointer-2}{}\fi
\noindent\textit{Source reference.} Original reference 07UY: sources/stacks-spaces-cohomology/pinned-original-source.tex, source lines 2991--2991. Navigation pointer for \textbackslash{}ref\{situation-vanishing\} used inside the counted span (line 3010); the argument that consumes it is retained in the span.\par\medskip
\ifdefined\hypertarget\hypertarget{sourcepointer-3}{}\fi
\noindent\textit{Source reference.} Original lemma 03M9: sources/stacks-chow/tags.txt, source lines 4741--4741. Navigation pointer for \textbackslash{}ref\{spaces-morphisms-lemma-pushforward\} used inside the counted span (line 3020); the argument that consumes it is retained in the span.\par\medskip
\begin{lemma}
\label{lemma-vanishing-compute}
In Situation \ref{situation-vanishing} for an $A$-module $M$ we have
$p_*(M \otimes_A \mathcal{O}_X) = \widetilde{M}$ and
$\Gamma(X, M \otimes_A \mathcal{O}_X) = M$.
\end{lemma}

\begin{proof}
The equality $p_*(M \otimes_A \mathcal{O}_X) = \widetilde{M}$ follows
from the equality $\Gamma(X, M \otimes_A \mathcal{O}_X) = M$ as
$p_*(M \otimes_A \mathcal{O}_X)$ is a quasi-coherent module on
$\Spec(A)$ by Morphisms of Spaces, Lemma
\ref{spaces-morphisms-lemma-pushforward}.
Observe that $\Gamma(X, \bigoplus_{i \in I} \mathcal{O}_X) =
\bigoplus_{i \in I} A$ by Lemma \ref{lemma-colimits}. Hence the
lemma holds for free modules. Choose a short exact sequence
$F_1 \to F_0 \to M$ where $F_0, F_1$ are free $A$-modules. Since
$H^1(X, -)$ is zero the global sections functor is right exact.
Moreover the pullback $p^*$ is right exact as well. Hence we see
that
$$
\Gamma(X, F_1 \otimes_A \mathcal{O}_X) \to
\Gamma(X, F_0 \otimes_A \mathcal{O}_X) \to
\Gamma(X, M \otimes_A \mathcal{O}_X) \to 0
$$
is exact. The result follows.
\end{proof}

\end{document}
